\newpage
\chapter{KESIMPULAN DAN SARAN} \label{Bab V}

\section{Kesimpulan} \label{V.Kesimpulan}
Berdasarkan hasil eksperimen dan analisis yang telah dilakukan terhadap metode HOUM pada tiga dataset dengan tingkat ketidakseimbangan yang berbeda, diperoleh kesimpulan sebagai berikut:
\begin{enumerate}
    \item Penerapan teknik \textit{oversampling} yang berbeda pada metode HOUM menghasilkan dampak yang beragam terhadap kinerja klasifikasi, sehingga tidak ditemukan satu teknik yang secara mutlak lebih unggul pada seluruh kondisi. Safe-Level SMOTE secara konsisten mampu meningkatkan deteksi kelas minoritas pada ketiga dataset karena teknik ini menempatkan data sintetis di area yang aman dan tidak tumpang tindih dengan kelas mayoritas. Di sisi lain, ADASYN menunjukkan keunggulan dalam menghasilkan keseimbangan klasifikasi yang lebih baik ketika diterapkan pada dataset yang memiliki jumlah fitur memadai dan batas antar kelas yang relatif jelas, sebagaimana terlihat pada dataset Diabetes. Akan tetapi, pada dataset yang memiliki fitur terbatas dan tingkat tumpang tindih yang tinggi seperti Haberman dan Transfusion, ADASYN justru menurunkan performa karena data sintetis yang dihasilkan berada di area yang sudah tercampur antar kelas. Oleh karena itu, pemilihan teknik \textit{oversampling} pada HOUM perlu disesuaikan dengan karakteristik dataset, terutama jumlah fitur dan tingkat tumpang tindih antar kelas.

    \item Variasi nilai parameter regularisasi ($C$) dan jenis kernel SVM terbukti berpengaruh terhadap kinerja metode HOUM. Nilai $C$ yang rendah cenderung menghasilkan proses \textit{undersampling} yang lebih efektif pada dataset dengan tingkat ketidakseimbangan yang tidak terlalu tinggi karena menghasilkan batas keputusan yang lebih umum. Sebaliknya, pada dataset dengan ketidakseimbangan yang tinggi, diperlukan nilai $C$ yang lebih besar untuk mencapai keseimbangan optimal antara deteksi kelas minoritas dan ketepatan prediksi. Meskipun demikian, nilai $C$ yang terlalu tinggi justru menyebabkan penghapusan data secara berlebihan sehingga menurunkan kinerja secara keseluruhan. Adapun pemilihan jenis kernel bergantung pada dimensi dataset, dimana kernel RBF lebih sesuai untuk dataset dengan fitur yang banyak karena mampu menangkap hubungan non-linear antar fitur, sedangkan kernel Linear sudah memadai untuk dataset berdimensi rendah. Kernel Polynomial tidak direkomendasikan karena cenderung mengalami \textit{overfitting} dan secara konsisten mencatat performa terendah di seluruh dataset yang diuji.
\end{enumerate}

\section{Saran} \label{V.Saran}
Berdasarkan hasil penelitian yang telah dilakukan, terdapat beberapa saran untuk pengembangan penelitian selanjutnya:
\begin{enumerate}
    \item Menggunakan dataset dengan dimensi fitur lebih tinggi dan kompleksitas lebih beragam untuk menguji konsistensi kernel RBF pada ruang fitur berdimensi tinggi.
    \item Memperluas rentang nilai $C$ ke bawah 0,5 untuk menganalisis pengaruh \textit{margin} yang lebih longgar terhadap proses \textit{undersampling}.
    \item Membandingkan teknik \textit{oversampling} lain seperti \textit{Borderline-SMOTE} dengan Safe-Level SMOTE dan ADASYN dalam kerangka HOUM.
    \item Menguji HOUM dengan algoritma klasifikasi lain seperti \textit{Gradient Boosting} atau \textit{Neural Network} untuk mengukur konsistensi peningkatan performa di berbagai model.
    \item Mengeksplorasi parameter tambahan SVM seperti \textit{gamma} pada kernel RBF untuk mengoptimalkan \textit{undersampling}.
    \item Menambah jumlah dan variasi rasio ketidakseimbangan dataset untuk mengevaluasi generalisasi HOUM secara lebih menyeluruh.
\end{enumerate}
